%!TEX root = ../main.tex
% =============================================================================
%  九、模型检验与灵敏度            建议 3–4 页
%  评审在问：你怎么证明是对的？
%  本节必须出现：① 三层验证（可行/最优/独立复算）
%                ② L1–L5 映射  ③ 消融矩阵（fig07）
%                ④ Sobol/单参数灵敏度 + \textbf{"结论是否改变"的判断}
%  红线：求解器返回 optimal 只说明"这个模型被解出来了"，不等于"模型建对了"。
%        任何"XX 是必须的"结论都必须有消融证据。
% =============================================================================

\section{模型检验与灵敏度分析}

\subsection{三层验证}

\begin{guideblock}
必须区分三层验证对象，并明确各层"不能证明什么"：
思路层（假设/信息结构/决策时序）不能靠数值结果反推；
模型层（变量/约束/单位/凸性）不能靠求解成功证明；
结果层（数值/口径/稳健性）不能靠一次最优解证明普适。
一句话模板见本节末，请照写并填入本队实测残差。
\end{guideblock}

验证对象分为三层：\textbf{思路层}（建模假设、信息结构、决策时序）不能靠数值结果反推；
\textbf{模型层}（变量、约束、目标、单位）不能靠"求解器返回最优"证明；
\textbf{结果层}（数值、口径、稳健性）不能靠一次最优解证明普适。
本文据此对三层分别设置检验动作，下面报告结果层的实测数值。

所有结果通过了原始可行性、对偶最优性与费用独立复算三重检验：
逐时段功率平衡的最大残差为 \ph~kWh（阈值 \(10^{-6}\)），
储电量递推的最大残差为 \ph~kWh，
对偶相对间隙为 \ph（阈值 \(10^{-6}\)），
费用独立复算与优化目标值的偏差为 \ph~元（阈值 \(10^{-6}\)）。

\subsection{五级验证链条}

\begin{table}[htbp]
  \centering
  \caption{验证链条 L1--L5 的实施方式、通过标准与本文结果}
  \label{tab:L1L5}
  \footnotesize
  \begin{tabularx}{\linewidth}{C{1.0cm} L{2.6cm} L{3.6cm} L{3.0cm} Y}
    \toprule
    级 & 名称 & 做法 & 通过标准 & 本文结果 \\
    \midrule
    L1 & 解析特例单测
      & 构造无光伏/无电池/单时段/两时段/常电价/\(\eta=1\) 等退化情形与闭式解对拍
      & 相对误差 \(\le 10^{-6}\)
      & \ph \\
    L2 & 合成数据端到端
      & 自造小规模数据算已知最优，让完整模型恢复；两阶段用"情景枚举 + 第一阶段枚举"对拍
      & 目标值一致、决策结构一致
      & \ph \\
    L3 & 真实数据独立审计
      & \textbf{不共用建模代码的费用函数}，只读原始附件与结果文件重算 10 项
      & 见表~\ref{tab:audit}
      & \ph \\
    L4 & 基线与消融
      & 四类基线 + 消融矩阵
      & 每个保留模块的移除造成可分辨的恶化（带区间）
      & \ph \\
    L5 & 统计与稳健性
      & 场景保真、样本外、配对比较、多种子 CI、灵敏度与盈亏平衡
      & 结论方向不变、关键量给区间、样本外无系统性反转
      & \ph \\
    \bottomrule
  \end{tabularx}
\end{table}

\subsection{独立审计项与阈值}

\begin{guideblock}
审计必须由\textbf{非实现者}执行：只从"输出文件"出发独立重算，
而不是读建模代码的逻辑——否则同一个符号约定错误会被复制一遍。
\textbf{失败即阻断}：任一 A1--A7 失败，禁止进入论文写作。
\end{guideblock}

\begin{table}[htbp]
  \centering
  \caption{独立审计项、阈值与本文实测结果}
  \label{tab:audit}
  \footnotesize
  \begin{tabularx}{\linewidth}{C{1.0cm} L{4.0cm} L{4.2cm} Y}
    \toprule
    项 & 检查内容 & 阈值 & 本文实测 \\
    \midrule
    A1 & 功率/电量平衡残差（式~\eqref{eq:q1-balance}）
       & \(\max_t\lvert E_t+PV_t\Delta t+d_t-L_t\Delta t-c_t-s_t\rvert\le 10^{-6}\)~kWh & \ph \\
    A2 & 储电量递推（式~\eqref{eq:q1-soc}）
       & \(\max_t\lvert S_{t}-S_{t-1}-\eta_cc_t+d_t/\eta_d\rvert\le 10^{-6}\)~kWh & \ph \\
    A3 & 储电量边界 & 越界数 \(=0\)（含边界值本身） & \ph \\
    A4 & 充放电功率边界 & \(c_t,d_t\le\overline P\Delta t\)，越界数 \(=0\) & \ph \\
    A5 & 充放互斥 & \(\min(c_t,d_t)\le 10^{-6}\) 全部成立 & \ph \\
    A6 & 周期/终端条件 & 问题 1：\(S_0-S_{144}=0\)；问题 2--4：跨日差 \(\le 10^{-6}\) & \ph \\
    A7 & 非负性 & 所有电量与费用 \(\ge -10^{-9}\) & \ph \\
    A8 & 费用复算 & 相对误差 \(\le 10^{-4}\)（计划 + 调整 + 违约 + 紧急），且差异可解释 & \ph \\
    A9 & 日期与结构 & 334 天；工作表名/列名/列序逐项一致 & \ph \\
    A10 & 全天汇总列 & 与逐时段求和一致（相对误差 \(\le 10^{-4}\)） & \ph \\
    \bottomrule
  \end{tabularx}
\end{table}

审计结论：\fbox{\ \ \ 全部通过，可进入写作\ \ \ }\quad/\quad%
\fbox{\ \ \ 存在失败项，已定位并修正\ \ \ }
\note{定稿前请只勾选一个，并把审计报告作为附录 D 附上}

\subsection{基线与消融}

\begin{guideblock}
消融的\textbf{公平性三条铁律}：同一场景集（先固定并存盘，所有变体读同一份）、
同一指标定义与口径、同一预算。
诚实检查：若某一层的边际贡献 \(<0.5\%\)，
\textbf{就写"该模块贡献有限，其主要价值在于保证可行性/提供解释，而非降低费用"}——
主动承认一个模块收益有限，比硬吹它有用可信得多。
\end{guideblock}

\begin{table}[htbp]
  \centering
  \caption{消融矩阵：八层基线（B0--B7）与完整模型（M）在同场景集上的费用对比}
  \label{tab:ablation}
  \footnotesize
  \begin{tabularx}{\linewidth}{C{1.0cm} L{4.6cm} L{2.8cm} C{2.4cm} Y}
    \toprule
    层级 & 模型配置 & 新增模块 & 总费用（元） & 相对 M 的差异 \\
    \midrule
    B0 & 无储能：全部净负荷按价购电 & ---（闭式，\(C^{base}\)） & \ph & \ph \\
    B1 & + 储能，但采用贪心策略 & 储能（无优化） & \ph & \ph \\
    B2 & + 储能，单日独立 LP（每日储电量清零） & 精确优化，但无跨日 & \ph & \ph \\
    B3 & + 跨日储电量连续递推 & 跨日状态 & \ph & \ph \\
    B4 & + 弃光变量 & 可行性保障 & \ph & \ph \\
    B5 & + 终端水价值（替代硬边界） & 防短视 & \ph & \ph \\
    B6 & + 不确定性（场景 + 两阶段） & 随机优化 & \ph & \ph \\
    B7 & + 多阶段滚动 + 非预期性 & 信息结构 & \ph & \ph \\
    M  & + 波动电价 + 联合情景 + 风险度量 & 风险决策 & \ph & 0 \\
    \bottomrule
  \end{tabularx}
\end{table}

\begin{figure}[htbp]
  \centering
  \includefig[width=0.95\linewidth]{fig07_ablation_matrix.pdf}
  \caption{消融矩阵（行：模型变体；列：总费用、紧急购电量、弃光电量、平均储电量利用率；
           单元格数值为该变体相对完整模型 M 的差异百分比；
           发散色标以 0 为中心；不可行单元格以灰色斜纹标出；
           全部变体在同一固定场景集上求解）}
  \label{fig:ablation}
\end{figure}

由表~\ref{tab:ablation} 与图~\ref{fig:ablation} 可得：
从 B0 到 B3 的改进贡献最大（合计 \ph），说明\textbf{跨日储电量递推}是本题的核心机制；
B4（弃光变量）对总费用的贡献仅 \ph，但其价值在于\textbf{保证极端日可行}；
B5--B7 与 M 的边际贡献分别为 \ph、\ph、\ph。

\subsection{灵敏度分析与结论反转点}

\begin{guideblock}
每个灵敏度都必须给出\textbf{"结论是否改变"}的判断，而不是只贴一张表。
禁止写"所有参数都重要，说明模型很敏感"。
龙卷风图\textbf{必须标出结论反转点}。
扰动必须作用到\textbf{运行区间}——只改物理最大容量而运行区间不变等于没有扰动。
\end{guideblock}

本文对表~\ref{tab:sens} 所列参数做单参数扫描与 Sobol 全局灵敏度分析，
关键量为总费用与储能搬运收益，稳定性用 20 个随机种子的均值 \(\pm\) 标准差与 95\% 置信区间刻画。

\begin{table}[htbp]
  \centering
  \caption{灵敏度分析：参数档位、结果区间与"结论是否改变"的判断}
  \label{tab:sens}
  \footnotesize
  \begin{tabularx}{\linewidth}{L{2.8cm} L{3.4cm} C{2.4cm} C{2.2cm} Y}
    \toprule
    参数 & 扰动档位 & 总费用区间（元） & Sobol 主指数 & 结论是否改变 \\
    \midrule
    充放电效率 \(\eta_c,\eta_d\) & \(0.85\)、\(0.90\)、\(0.95\) + 往返 81\% 口径 & \ph & \ph & \ph \\
    储电量运行区间 \([\underline S,\overline S]\) & \([600,11400]\)、\([1200,10800]\)、\([2400,9600]\) & \ph & \ph & \ph \\
    充放功率上限 \(\overline P\) & 2500、5000、8000~kW & \ph & \ph & \ph \\
    初始储电量 \(S^{init}\) & 1200、6000、10800~kWh & \ph & \ph & \ph \\
    紧急购电倍数 \(\kappa\) & 2、3、5、10 & \ph & \ph & \ph \\
    调整倍数 \((\kappa^+,\kappa^-)\) & \((1.2,0.2)\)、\((1.5,0.5)\)、\((2.0,0.8)\) & \ph & \ph & \ph \\
    终端条件 & 自由 / \(S_T\ge S_0\) / 日周期 & \ph & \ph & \ph \\
    电价水平 & 附件 4 \(\pm20\%\)、峰谷差压缩 50\% & \ph & \ph & \ph \\
    预报误差水平 & \(\times0.5/\times1/\times2\) & \ph & \ph & \ph \\
    负载预测误差 & 相对 MAE 0/5\%/10\% & \ph & \ph & \ph \\
    \bottomrule
  \end{tabularx}
\end{table}

\begin{guideblock}
"结论是否改变"一列必须写出\textbf{具体的反转条件}，例如：
"当 \(\kappa<2\) 时，报童临界分位数降至 \ph 以下，'\(0{:}00\) 一次决策已足够'的结论不再成立"。
若某参数的扫描区间内结论始终不变，就写"该区间内结论不变"，这也是有效结论。
\end{guideblock}

\subsection{样本外回测与验证话术自查}

本文的验证协议在建模前锁定：以 2025-01-01 至 2025-09-30 为拟合窗口
（用于场景分布拟合与参数选择），以 2025-10-01 至 2025-12-31 为测试窗口；
\textbf{测试窗口未参与任何调参}，仅用于最终报告。
测试窗口上的总费用为 \ph~元，相对拟合窗口的平均水平偏差 \ph，
未出现系统性反转。2025 年 1 月作为储电量递推的预热期，不计入结果（假设 A8）。

本文不使用被禁止的验证话术。八项伪验证自查（自证循环、以求解器状态代替物理校验、
比较不同样本、以训练场景充当样本外、只报点估计、事后挑选指标、以"结果好看"代替"模型正确"、
只改物理容量而不改运行区间）已在审计脚本中以断言形式固化，逐条结果见附录 D。
